% problem-000345 9 有机骨架与杂化分析
\section*{9. 有机骨架与杂化分析}
（图：）

\subsection*{9-1}
(1)的反应条件是什么？(1)的反应类别是什么？(2)的反应类别是什么？

\subsection*{9-2}
分⼦ A 中有 个⼀级碳原⼦，有 个⼆级碳原⼦，有 个三级碳原⼦，有 个四级碳原⼦，⾄少有个氢原⼦共平⾯。

\subsection*{9-3}
B的同分异构体D的结构简式是：

\subsection*{9-4}
E是A的⼀种同分异构体，E含有sp、sp2、sp3杂化的碳原⼦，分⼦中没有甲基，E的结构简式是：第10题（15分）⾼效低毒杀⾍剂氯菊酯(I)可通过下列合成路线制备：

（图：）

（图：）

（图：）

\subsection*{10-1}
化合物A能使溴的四氯化碳溶液褪⾊且不存在⼏何异构体。画出A和B的结构简式。

\subsection*{10-2}
化合物E的系统名称是 ，化合物I中官能团的名称是 o

\subsection*{10-3}
由化合物E⽣成化合物F经历了 步反应。每步反应的反应类别分别是 C

\subsection*{10-4}
在化合物E转化成化合物F的反应中，能否⽤ $\mathrm { N a O H / C _ { 2 } H _ { 5 } O H }$ 代替 $\mathrm { C _ { 2 } H _ { 5 } O N a / C _ { 2 } H _ { 5 } O H }$ 溶液？为什么？

\subsection*{10-5}
化合物G和H反应⽣成化合物I、 $\mathrm { { N } ( C H _ { 2 } C H _ { 3 } ) _ { 3 } }$ 和 $\mathrm { N a C l }$ ，由此可推断：H 的结构简式是 ，H 分⼦中氧原⼦⾄少与 个原⼦共平⾯。

\subsection*{10-6}
芳⾹化合物J⽐F少两个氢，J中有三种不同化学环境的氢，它们的数⽬⽐是9:2:1，则J可能的结构为（⽤结构简式表⽰）：
